\chapter{Exploratory Cross-Dataset Check on ADS-24}
\label{ch:ads24}

Because the in-house results of Chapter~\ref{ch:results} come from one
laboratory, we repeated the study on data collected elsewhere. We
did not find a public dataset that could reproduce the PROPHET task as
defined in Chapter~\ref{ch:task}, so we relaxed it: alerts ordered by
their occurrence time stand in for the receipt boundary of
Section~\ref{sec:receipt}, the dataset's attack-stage labels for
source-designated attacker actions, and a held-out condition with
simulated users for a held-out benign worker. Under these relaxations
we used ADS-24, a dataset of attack episodes in a cyber range with
simulated user agents~\cite{nyberg2026cyber}; in the release we
inspected, each episode records whether the users were active.\footnote{KTH
data repository, DOI 10.71775/kth.4mh6z-zw065, distributed under CC
BY-SA 4.0 according to the record metadata; accessed 20 November 2026.}
% 원문(nyberg2026cyber): arXiv 2609.16541v2 초록 "a virtualized computer network with variable topologies, red-team automation
%   and simulated user agents"; V절 a) Simulated Users(사용자 있음·없음 실행). 확인됨.

This chapter is therefore an exploratory check on a derived task. The
132 runs are ones we built from the downloaded files, one per episode,
with labels from a documented rule binding attack stages to alert
times, which we did not rerun (Appendix~\ref{app:ads24-data}, which
also gives why six other candidates were not used). The results are
development results and show whether the in-house direction appears at
all on data collected elsewhere.
% 5회차 방법 리뷰 P1-3: ADS-24를 탐색적 확인으로 낮췄다. 장 위치는 decisions.md 37번(본문 유지).
% 6회차 문장 P2-10, aitone P3-2: 긴 문장을 나누고 범위를 긍정형으로 먼저 적었다.
% 출처(배포·사용권·접근일): sources/atlas_answers/ads.md 28번; atlas
%   OT/PROPHET_ENV_CONSENSUS_KD_V3_20261120/ADS24_V3_DRAFT.md:5,17,18 (KTH, DOI, CC BY-SA 4.0, 2026-11-20).
%   사용권은 license 파일이 아니라 기록 API의 metadata.rights에서 읽은 것이다.
% 사용자 on/off: 인용 리뷰 P2-5에 따라 우리가 검사한 배포본 기준으로 적었다.
% 근거: RECORD:212, 236, 265--275 (선정 경위, 후보 7행 표는 부록 D).
% ADS24_CONTRACT_SEALED.md "과제와 자료": cutoff 이전 Wazuh @timestamp 경보만
% 입력. 이 시각이 수신 시각인지 확인되지 않아 계약은 "경보 시각 예측"이라 부른다.


\section{Task and data roles}
\label{sec:ads24-task}

At each prediction boundary the input is the sequence of Wazuh alerts
recorded before it, read through the Wazuh rule identifier and rule
level, \texttt{rule.id} and
\texttt{rule.level},\footnote{Wazuh rule syntax:
\url{https://documentation.wazuh.com/current/user-manual/ruleset/ruleset-xml-syntax/rules.html},
accessed 24 November 2026.}
together with host role and alert timing. The target is the ATT\&CK tactic of the next attack segment, one
of nine classes. A segment merges consecutive stages that share a
tactic, and every segment after the first starts a boundary. The rule was fixed before any model was scored. These are oracle
boundaries, placed where the labeled tactic sequence changes, so every
result in this chapter is conditional next-tactic classification at
those transitions.

The contract assigns the 132 runs to four roles. Without simulated
users, the seen condition, 48 runs form FIT (8,461 boundaries) and 15
check preservation (2,982). With users, the unseen condition, 34 runs
supply unlabeled adaptation inputs (5,811) and 35 others form the
environment evaluation (5,639), which the CD student never reads. The runs without users were
collected mostly in early April and those with users in mid-April, so
the results do not estimate an effect of the simulated users alone.
% 2026-11-25 분량 줄이기 C13: 역할 manifest·실행별 감사 파일 문장은 아래 출처 주석과 부록 E 의 manifest 목록에 맡겼다.
% 132 대 134, episode=run, 시각 결속 규칙은 부록 D에 적었다(ads.md 25·26번).
% 출처(역할 manifest): 원래 PC: PROPHET_ENV_CONSENSUS_KD_V3_20261120/ADS24_STAGE1/ROLE_MANIFEST.json:583
%   (roles: FIT_OFF 48, PRESERVE_OFF 15, ADAPT_ON 34, EVAL_ON 35 — 본문 48/15/34/35와 같다).
% 출처(실행별 감사): 원래 PC: PROPHET_ENV_CONSENSUS_KD_V3_20261120/ADS24_EPISODE_TIME_AUDIT.csv:1
%   (label_exclusion, time_exclusion, error 열).
% 근거: RECORD:214, 247, 251, 254; ADS24_CONTRACT_SEALED.md "과제와 자료",
% "입력과 모델 이식"; ADS24_STAGE2_REPORT.md "계보" (8,461/2,982/5,811/5,639경계).
% 연속 경계의 약 62.7%가 같은 전술이라는 수치는 개요에 따라 부록 D에 둔다.
% 원 보고서의 REL_BASE = 이 장의 REL_GLOBAL-ADS, KD_ENV_TOP5 = CD student,
% KD_SELF = self-distillation control.
% RECORD:441 (OFF 4월 초, ON 4월 중순). Wazuh 문서는 answers 07:10 2번에 따라
% 각주 URL로 바꿨다. 수집에 쓰인 Wazuh 버전은 우리가 검사한 KTH 기록에 적혀 있지 않다(ads.md 31번).
% (해결, 저자 확인 31번) Wazuh 문서 URL의 접근일은 2026-11-24 에 열어 확인한 날로 각주에 적었다.

Two stages were run, each under its own sealed contract. Stage~1
trained 30 REL\_GLOBAL-ADS seeds (Section~\ref{sec:method-transfer})
and passed its gate, a seed standard
deviation of at least 0.050 on the environment evaluation, with 0.163
$[0.140, 0.191]$. Stage~2 then trained 60 new models from the same 30
seeds: a CD student whose teachers are the five REL\_GLOBAL-ADS seeds with
the highest top-1 agreement on the 34 adaptation runs, and a
self-distillation control, with all other settings unchanged. Because
stage~1 had already scored the 35 environment-evaluation runs and used
them to decide that stage~2 should run, stage~2 is a development
comparison.
% 근거: RECORD:256--261; ADS24_CONTRACT_SEALED.md "내부 선택과 1단계" (관문 SD≥.050,
% 10,000회 95% 구간 하단≥.020), "2단계 예정 계약" (교사: ADAPT_ON34에서만,
% 자기 포함 30모델 다수결; KD_SELF는 같은 seed의 동결 REL_BASE);
% ADS24_STAGE1_REPORT.md:3 (.162925 [.140425, .190702]); ADS24_STAGE2_REPORT.md
% "계보" (새 학생 60 fit, 교사 301030, 301026, 301025, 301003, 301022, KL 교사‖학생,
% 온도 1, 계수 .25). 1단계의 별도 pilot 9 fit은 표 tab:rounds에 있다.

\section{Results}
\label{sec:ads24-results}

All three pre-specified contrasts fell on the favorable side of their
thresholds (Tables~\ref{tab:ads24-groups} and~\ref{tab:ads24-criteria}). On the
environment evaluation the seed standard deviation fell from 0.163 to
0.025 and the mean NLL from 2.276 to 2.060, and the preservation
interval stays well below the limit of $+0.043$. The self-distillation
control left the environment loss where it was; as a secondary contrast
with individual 95\% intervals, the CD student minus the control was
$-0.215$ $[-0.248, -0.181]$ in mean NLL and $-0.137$
$[-0.163, -0.115]$ in seed standard deviation.
% 근거: RECORD:216--227, 232; ADS24_STAGE2_REPORT.md "세 군의 전체 성과",
% "봉인한 주 판정", "TOP5−SELF 보조 대비" (env_mean −0.214832 [−0.248399, −0.181309];
% env_sd −0.136728 [−0.162917, −0.114904]). 해결된 label TODO는 지웠다.

\begin{table}[htbp]
\centering
\caption[The three groups of 30 seeds on ADS-24]{The three groups of 30 seeds on ADS-24. Environment columns
are the 35 runs with simulated users (5,639 boundaries); preservation is
the 15 runs without users (2,982 boundaries). NLL is execution-weighted
as in Equation~\ref{eq:nll}; accuracy is top-1 accuracy. The last
column counts the seeds whose environment NLL is lower than that of the
FIT prior defined below.}
\label{tab:ads24-groups}
\small
\setlength{\tabcolsep}{4pt}
\begin{tabular}{@{}lrrrrrr@{}}
\toprule
Group & Env.\ NLL & Seed SD & Worst seed & Env.\ accuracy
  & Pres.\ NLL & Seeds below FIT prior \\
\midrule
REL\_GLOBAL-ADS            & 2.276 & 0.163 & 2.690 & 28.4\% & 1.449 & 0/30 \\
CD student                 & 2.060 & 0.025 & 2.121 & 31.5\% & 1.430 & 0/30 \\
Self-distillation control  & 2.275 & 0.162 & 2.681 & 28.2\% & 1.438 & 0/30 \\
FIT prior                  & 1.830 & ---   & ---   & 33.9\% & ---   & --- \\
\bottomrule
\end{tabular}
\end{table}
% 근거: RECORD:216--221; ADS24_STAGE2_REPORT.md "세 군의 전체 성과", "FIT 다수 전술
% 기준선과 비교" (0/30 세 군 모두). 정확도는 원 파일의 0.283592, 0.315363, 0.282060,
% 0.339311을 소수 첫째 자리 %로 줄였다.

\begin{table}[htbp]
\centering
\caption[Pre-specified contrasts of the ADS-24 second stage]{Pre-specified contrasts of the ADS-24 second stage, a
development comparison: CD student minus REL\_GLOBAL-ADS, with two-sided
Bonferroni-adjusted 98.3333\% intervals from 100,000 execution resamples
stratified by scenario. The 30 seeds are fixed. The 35
environment-evaluation runs had been scored in stage~1.}
\label{tab:ads24-criteria}
\small
\begin{tabular}{@{}llrlrl@{}}
\toprule
Runs & Contrast & Estimate & Interval & Limit & Outcome \\
\midrule
Environment  & seed SD difference & $-0.138$ & $[-0.167, -0.113]$ & $0$
  & favorable \\
Environment  & mean NLL difference & $-0.216$ & $[-0.257, -0.174]$ & $0$
  & favorable \\
Preservation & mean NLL difference & $-0.019$ & $[-0.024, -0.015]$
  & $+0.043$ & favorable \\
\bottomrule
\end{tabular}
\end{table}
% 근거: RECORD:223--227; ADS24_STAGE2_REPORT.md "봉인한 주 판정" (−0.137581
% [−0.167053, −0.113060]; −0.215652 [−0.256788, −0.174157]; −0.019398
% [−0.024140, −0.015072]; 한도 +0.043482, PCG64 2026112003).

As in-house, the seeds spread far more on the unseen condition (0.163
against 0.015 on the preservation runs), and the seen condition says
little about which seed will hold up there: post hoc, the Spearman
correlation between the seeds' preservation and environment losses was
about $-0.39$, against $-0.131$ in-house. The in-house agreement
statistic S1 was not computed; a stage-1 diagnostic instead correlated
each seed's top-1 agreement with the majority of the other 29 models
with its NLL, expecting $\rho \le -0.50$, and gave $\rho = +0.127$
(permutation $p = 0.503$).
% 근거: RECORD:231; ADS24_STAGE2_REPORT.md (REL_BASE 보존 seed SD 0.015182).
% 출처(−0.39): atlas OT/PROPHET_ENV_CONSENSUS_KD_V3_20261120/ADS24_STAGE1_REPORT.md:23-54
%   (30 seed 표: 열 2 환경 NLL, 열 4 보존 OFF NLL). 원 보고서 수치가 아니라 이 표에서 다시 계산한
%   값이다: ρ=−0.38776 (n=30, NLL 대 NLL, 동순위 없음), 보존 SD 0.015182 (diag.md 끝, 저자 답 4번).

% 근거: RECORD:233, 260, 405; ADS24_CONTRACT_SEALED.md "내부 선택과 1단계" (Spearman
% ρ≤−.50, 100,000회 양측 순열 p<.05; 진단 조건이며 증류 관문이 아님);
% ADS24_STAGE1_REPORT.md:15 (ρ .127475, p .503125); 30 seed 표(:25-54)에서 합의도(6열)–환경
%   NLL(2열) Spearman을 다시 계산해 0.127475 확인(ads.md 39번).
% 6회차 방법 리뷰 P2-2: "not reproduced"를 사전 지정 진단의 불리한 결과로 한정했다.

Absolute performance stayed low. The \emph{FIT prior}, the
run-weighted tactic frequencies of the 48 FIT runs with add-one
smoothing, scored 1.830, and all 30 CD seeds were worse.
% 출처: atlas OT/ARCH_GENERALITY_STAGE1_20261122/ARCH_GENERALITY_STAGE1_REPORT.md:71,72 (ads.md 27번).
% 출처(+1 평활): 원래 PC: PROPHET_ENV_CONSENSUS_KD_V3_20261120/ads24_stage1.py:574
%   (p=(counts+1)/(counts.sum()+len(LABELS))). 옛 원고의 "원장에 평활 표시 없음"은 코드와 달라 고쳤다.
The CD mean exceeded it by $+0.231$, against $+0.446$ for
REL\_GLOBAL-ADS, and no CD seed reached the best REL\_GLOBAL-ADS seed
(1.953).

Two later analyses outside the contract point the same way. For the
plain GRU, students of five random teachers did better on ADS-24 than
students of the five most agreeing ones (Section~\ref{sec:gen-teachers}).
And on the 119 usable runs outside the contract's 132, which carry the
flags the contract excluded and so form another population, the frozen
REL\_GLOBAL-ADS models moved from 1.948 to 1.862 in NLL and from 0.069
to 0.019 in seed standard deviation, and the three general encoders
moved the same way (Appendix~\ref{app:ads24-data}).
% 출처: inputs/data/FINAL_CHECKS_20261125/UNIFIED_REANALYSIS_20261125/ADS_UNUSED_119_GROUPS.csv (REL_GLOBAL BASE 1.948317/0.069278,
%   KD_ENV 1.862240/0.018549; LSTM 1.949566→1.878155, LogBERT 1.939096→1.788063, PLAIN_GRU 1.937812→1.903430);
%   UNIFIED_REANALYSIS.md:67--69 (사후, 계약 제외 모집단, 플래그 겹침).
% 2026-11-25 C13: "이동에 따른 손실은 CD 뒤에도 남는다"는 12.1절(sec:disc-does)에서 다룬다.
% 근거: RECORD:212, 221, 234, 423; ADS24_STAGE2_REPORT.md "FIT 다수 전술 기준선과
% 비교" (기준선 1.829580; 평균 차 +0.446462, +0.230810, +0.445642; 0/30) 및
% 최고 seed (REL_BASE 301008 1.953153, KD_ENV_TOP5 301008 2.020009).
% 이 기준선(FIT만으로 정한 다수 전술, 환경 ON35)은 부록 D의 라벨 기준선
% (132실행 5-fold, 1.8547/32.0%)과 다르다.
% 정의: METHOD_EXTENSION 원장의 "실행 동등 FIT 사전 빈도"(FIT 다수 전술 1.829580).

% 2026-11-25 분량 줄이기 C13: 사전 지정 보조 분해와 사후 교사 평균 두 문단을 부록 D 로 옮겼다.
% 개요는 이 장을 7문단으로 잡았으나, 지위(1·2단계, 열람 여부)를 과제 문단과 나눠
% 8문단이 되었다.
